\section*{Bab \ref{ch:b3:radicals-carbenes} --- Radikal, Karbena, dan Penataan Ulang}

\begin{solution}{exo:b3:radicals-carbenes:1}
Propana memiliki enam hidrogen primer dan dua hidrogen sekunder. Klorinasi: $6 \times 1 : 2 \times 3.9 = 6 : 7.8$,
yaitu 43\,\% 1-kloropropana dan 57\,\% 2-kloropropana. Brominasi: $6 : 164$, 4\,\% dan
96\,\%.
\end{solution}

\begin{solution}{exo:b3:radicals-carbenes:2}
$\ce{CH2}$: triplet; $\ce{CCl2}$: singlet (pasangan elektron bebas klor menaikkan orbital p yang kosong).
Diklorokarbena teradisi secara stereospesifik: \textit{trans}-1,1-dikloro-2,3-dimetilsiklopropana.
\end{solution}

\begin{solution}{exo:b3:radicals-carbenes:3}
\textit{cis}-1,2-Dietilsiklopropana: \omterm{def:b3:radicals-carbenes:carbenoid}{karbenoid} teradisi dalam satu tahap pada satu muka.
\end{solution}

\begin{solution}{exo:b3:radicals-carbenes:4}
Sikloheksanon: $\varepsilon$-kaprolakton (oksepan-2-on). Asetofenon: fenil asetat,
$\ce{CH3CO2C6H5}$; fenil lebih dahulu bermigrasi daripada metil.
\end{solution}

\begin{solution}{exo:b3:radicals-carbenes:5}
Sembilan hidrogen primer dan satu hidrogen tersier. Klorinasi: $9 : 5.1$, 64\,\% 1-kloro-2-metilpropana
dan 36\,\% 2-kloro-2-metilpropana. Brominasi: $9 : 1600$, 0.6\,\% dan 99.4\,\%.
\end{solution}

\begin{solution}{exo:b3:radicals-carbenes:6}
$k_c = 3.0 \times k_H[\mathrm{Bu_3SnH}] = 3.0 \times \num{2.4e6} \times 0.020 = \qty{1.4e5}{s^{-1}}$.
\end{solution}

\begin{solution}{exo:b3:radicals-carbenes:7}
Serangan pada C5: 5-ekso-trig, disukai; pada C6: 6-endo-trig, disukai menurut aturan itu tetapi
lebih lambat di sini. 4-ekso-tet: disukai. 5-endo-trig: tidak disukai. 5-ekso-dig: disukai.
3-ekso-dig: tidak disukai. 5-endo-dig: disukai.
\end{solution}

\begin{solution}{exo:b3:radicals-carbenes:8}
Lepasnya air memberikan kation tersier yang benzilik; pada karbon tetangga, gugus fenil
dan gugus metil dapat berpindah, dan fenil, yang \omterm{def:b3:radicals-carbenes:migratory}{kecenderungan migrasinya} lebih tinggi, bermigrasi. Kation
yang distabilkan oksigen itu melepaskan sebuah proton: 3,3-difenilbutan-2-on.
\end{solution}

\begin{solution}{exo:b3:radicals-carbenes:9}
Curtius: $\ce{C6H5CON3 -> C6H5NCO + N2}$. Dengan air, asam karbamat melepaskan $\ce{CO2}$: anilina
(sebagian difenilurea terbentuk ketika anilina bertemu dengan \omterm{def:b3:radicals-carbenes:curtius}{isosianat} lagi). Dengan etanol: etil
\textit{N}-fenilkarbamat, $\ce{C6H5NHCO2C2H5}$.
\end{solution}

\begin{solution}{exo:b3:radicals-carbenes:10}
Klor: primer $421.7 - 432 = \qty{-10}{kJ/mol}$, tersier $405.4 - 432 = \qty{-27}{kJ/mol}$; brom: $+56$ dan
$+39$. Kedua pengambilan oleh klor bersifat eksotermik, dengan keadaan transisi awal yang
hanya sedikit merasakan perbedaan \qty{16}{kJ/mol}; kedua pengambilan oleh brom bersifat endotermik,
dengan keadaan transisi akhir yang energi aktivasinya berbeda sebesar sebagian besar perbedaan itu:
$\eu^{16000/(8.314 \times 298)} \approx 600$, yaitu orde rasio yang teramati.
\end{solution}

\begin{solution}{exo:b3:radicals-carbenes:11}
Oksim \textit{E} (OH anti terhadap C2, yang membawa gugus metil): C2 bermigrasi, dan nitrogen masuk
di sebelahnya: 7-metilazepan-2-on. Oksim \textit{Z}: C6 bermigrasi: 3-metilazepan-2-on. Pada
\omterm{def:b3:radicals-carbenes:beckmann}{penataan ulang Beckmann}, geometri oksim yang menentukan; pada \omterm{def:b3:radicals-carbenes:baeyer}{oksidasi
Baeyer--Villiger}, \omterm{def:b3:radicals-carbenes:migratory}{kecenderungan migrasi} yang menentukan, dan karbon sekunder yang berpindah:
7-metiloksepan-2-on.
\end{solution}

\begin{solution}{exo:b3:radicals-carbenes:12}
Fraksi tersiklisasi $k_c/(k_c + k_H c)$: 0.91, 0.49, dan 0.09. Untuk 95\,\%: $k_Hc = k_c(1/0.95 - 1)$, sehingga
$c = 0.0526 \times \num{2.3e5}/\num{2.4e6} = \qty{5.0e-3}{mol/L}$. Konsentrasi itu dipertahankan dengan menambahkan hidrida perlahan-lahan dengan
pompa suntik, atau dengan memakai sejumlah katalitik timah halida bersama zat pereduksi
yang meregenerasi hidrida itu.
\end{solution}

\begin{solution}{pb:b3:radicals-carbenes:1}
\textbf{1.} $\ce{C6H10O + NH2OH -> C6H11NO + H2O}$.

\textbf{2.} Asam diperlukan untuk menyingkirkan gugus hidroksi zat antara adisi sebagai
air, tetapi asam yang terlalu banyak memprotonasi hidroksilamina, yang lalu tidak lagi teradisi.

\textbf{3.} Tidak: kedua karbon yang terikat pada karbon C=N setara karena
simetri cincinnya.

\textbf{4.} Pada oksim \textit{E}, OH mengarah menjauhi karbon C2 yang membawa metil (anti terhadap
karbon itu), sedangkan pada oksim \textit{Z} mengarah ke karbon itu.

\textbf{5.} Inversi pada nitrogen oksim yang membawa oksigen elektronegatif memiliki penghalang
yang tinggi.

\textbf{6.} Asam memprotonasi (atau menyulfonasi) gugus hidroksi, sehingga menjadikannya gugus pergi
yang baik, yaitu air.

\textbf{7.} Ikatan C--C cincin yang anti terhadap gugus pergi berpindah dari karbon ke nitrogen,
yang dengan demikian tersisip ke dalam cincin: enam atom menjadi tujuh.

\textbf{8.} Ion nitrilium siklik.

\textbf{9.} $\ce{C6H11NO -> C6H11NO}$: sebuah isomerisasi, dari oksim menjadi laktam.

\textbf{10.} Kaprolaktam, azepan-2-on, yaitu amida siklik beranggota tujuh dari
asam 6-aminoheksanoat; polimerisasi pembukaan cincinnya memberikan nilon-6.

\textbf{11.} \qty{98.0}{g/mol} dan \qty{113.0}{g/mol}.

\textbf{12.} $\qty{1.0e6}{g}/\qty{98.0}{g/mol} = \qty{10204}{mol}$; $\times 0.98 \times 0.98 = \qty{9800}{mol}$; $\times \qty{113.0}{g/mol} = \qty{1.11e6}{g}$, \qty{1.11}{t}.

\textbf{13.} Oksim: $\num{10204} \times 0.98 = \qty{1.00e4}{mol}$; $\ce{H2SO4}$: $\qty{1.30e4}{mol}$, yang memberikan jumlah mol
$\ce{(NH4)2SO4}$ yang sama: $\qty{1.30e4}{mol} \times \qty{132.1}{g/mol} = \qty{1.7}{t}$, lebih banyak daripada kaprolaktamnya.

\textbf{14.} $\ce{C6H10O + RCO3H -> C6H10O2 + RCOOH}$.

\textbf{15.} Hasil adisi tetrahedral asam peroksi pada karbon karbonil (zat antara
Criegee), sebuah peroksi hemiketal.

\textbf{16.} Sebuah gugus $\ce{CH2}$ cincin (keduanya setara) berpindah ke oksigen peroksida
yang lebih dekat; karboksilat $\ce{RCOO-}$ lepas.

\textbf{17.} $\varepsilon$-Kaprolakton, oksepan-2-on; polimerisasi pembukaan cincinnya
menghasilkan poliester yang terbiodegradasi, yaitu poli($\varepsilon$-kaprolakton).

\textbf{18.} Satu-satunya produk sampingnya adalah air, sedangkan asam peroksi meninggalkan satu mol
asam karboksilat per mol lakton dan penyimpanannya sendiri berbahaya.

\textbf{19.} Karbon sekunder C2 bermigrasi: 7-metiloksepan-2-on.

\textbf{20.} Ya: karbon yang bermigrasi mempertahankan konfigurasinya.

\textbf{21.} 7-Metilazepan-2-on (C2 anti, bermigrasi).

\textbf{22.} 3-Metilazepan-2-on (C6 bermigrasi).

\textbf{23.} Beckmann: geometri oksim (gugus yang anti berpindah).
Baeyer--Villiger: \omterm{def:b3:radicals-carbenes:migratory}{kecenderungan migrasi} (karbon yang lebih tersubstitusi berpindah).

\textbf{24.} Kedua karbon $\alpha$-nya setara: satu oksim, satu laktam, satu lakton.

\textbf{25.} Satu ton sikloheksanon memberikan sekitar \qty{1.11}{t} kaprolaktam.
\end{solution}
